% ── Part III, Chapter 3 ───────────────────────────────────────
% Spectra — learning to read stellar light as encoded data,
% inference through transformation
% Payne-Gaposchkin, Harvard, 1923--1924

\chapter{Spectra}

\strandmarker{Payne-Gaposchkin}{Harvard, 1923--1924}

\lettrine{T}{he} plates arrive from the night before.

Someone has been at the telescope while she slept,
tracking the motion of the sky, holding the instrument
steady against the drift, collecting the light that
has been travelling, in some cases, for longer than
the human species has existed. The light arrives and
strikes the photographic emulsion and leaves a mark.
The mark is not the light. It is what the light did
to the material it encountered, which is a different
thing, and which is all that is available.

She holds a plate to the light.

The spectrum is a band of colour interrupted by lines:
dark gaps where specific wavelengths are absent,
absorbed by atoms in the stellar atmosphere that took
exactly that energy and left nothing to continue the
journey. Each line is the signature of an element in
a specific state. Each line's position is fixed by
physics. Each line's darkness---its depth, its width---
encodes something about how much of that element is
present and under what conditions.

The plate is a message.

She is learning the language.

\section{}

The language has a grammar.

It is not a grammar she invented or that anyone fully
invented: it emerged from the accumulated work of
people who looked at spectra and noticed regularities,
who noticed that the regularities corresponded to
something in the structure of atoms, who built the
correspondence into a theory that could be run in
both directions---from atom to spectrum, or from
spectrum back to atom.

Running it backward is the harder direction.

Forward is deduction: given the atom, predict the
spectrum. This is clean, constrained, deterministic
within the theory's domain. Backward is inference:
given the spectrum, reconstruct the atom, or more
precisely, reconstruct the conditions under which
the atom produced this spectrum. The backward
direction is harder because multiple conditions
can produce similar spectra, because the plate
is imperfect, because the atmosphere between the
star and the instrument introduces its own
signatures that must be subtracted before the
star's own signal can be read.

She learns to subtract the atmosphere.

She learns to distinguish the instrument's
characteristic distortions from the signal the
instrument is carrying. She learns, in other
words, to read through the layers of transformation
that stand between the event---the atom absorbing
the photon in the stellar atmosphere---and the
record---the dark line on the glass plate in
her hands.

This is not simple. It requires holding multiple
frameworks simultaneously: the physics of atomic
transitions, the physics of stellar atmospheres,
the optics of the telescope and spectrograph,
the chemistry of the photographic emulsion. Each
framework constrains the interpretation. Where
the constraints converge, the interpretation is
strong. Where they conflict, something is wrong,
and the wrongness is information about which
framework requires revision.

\section{}

She develops a method, though she is careful
about that word.

What she develops is more like a discipline of
attention: a way of approaching the plates that
does not impose an interpretation before the
plates have been allowed to constrain one. She
has seen, in the literature, what happens when
the interpretation comes first---when a framework
is brought to the data with enough confidence
that the data are read as confirming it, the
anomalies smoothed over or set aside for later
attention that never arrives.

She does not bring confidence.

She brings frameworks, plural, held loosely,
each one a lens that reveals certain features
while obscuring others. She moves between them,
applying each in turn, noting where they agree
and where they diverge, treating the divergences
not as problems to be resolved by choosing a
winner but as the most informative regions of
the data---the places where the evidence is
putting the most pressure on the assumptions.

The plates accumulate.

She is working with more stellar spectra than
have ever been systematically analysed by a
single researcher, a volume of data that creates
its own kind of constraint: patterns that appear
in a single spectrum might be artifacts, but
patterns that appear across hundreds of spectra,
in stars of different types at different distances,
are telling her something about the stars rather
than about her instruments or her methods.

She begins to see the pattern.

\section{}

The pattern concerns hydrogen.

It is everywhere. Not as a trace, not as one
element among many in roughly equal proportions,
but as a presence so dominant that the other
elements, taken together, do not approach it.
The spectra say this clearly, once the theory
of ionisation is applied correctly---once she
accounts for the fact that most hydrogen in
stellar atmospheres is in states that do not
produce visible absorption lines, that the
absence of lines is not absence of the element.

This is the move that the previous analyses
had not made.

They had looked at the lines and inferred the
composition from the lines, without fully
accounting for what the absence of lines could
mean. The inference was locally reasonable---
the lines you see correspond to the elements
present---but globally incomplete, because the
lines you don't see also correspond to elements
present, in states the lines don't reveal.

She applies the ionisation theory and the
data reorganise themselves.

The reorganisation is not subtle. It is not
a small correction to a previously adequate
picture. It is a restructuring of the picture
entirely, in which what was taken as background
becomes foreground, what was treated as minor
becomes dominant, and the composition of
stars---which has been assumed to resemble
the composition of the earth because the
earth was the composition that was familiar
---turns out to be something else entirely.

\section{}

She checks the derivation.

Then she checks it again, because a result
this large carries the obligation of extra
scrutiny---not from doubt, exactly, but from
the recognition that the distance between
a correct result and a systematic error can
be difficult to see from inside the analysis
that produced either one.

The derivation holds.

She follows the chain of inference from the
plates to the ionisation theory to the line
strengths to the abundances, and at each link
she asks whether something else could produce
the same output, whether the constraint is
as tight as it appears, whether she has
accounted for what she thinks she has accounted
for.

The constraint is tight.

The result is not one interpretation among
several plausible ones. It is the interpretation
that the data, the physics, and the theory
jointly produce when none of them is overridden
in favour of prior expectation.

She writes it up.

The writing is careful and exact and does not
hedge where the evidence does not require
hedging. The conclusion is stated as what it
is: the conclusion the evidence warrants,
derived by the method the problem requires,
checked against the constraints available.

She does not know yet what will happen to it.

She knows what it says.

The stars are mostly hydrogen.

The universe is not what anyone thought it was.

